\documentclass[12pt]{article}

\begin{document}
Hello! This is my first \LaTeX\ document.\\
A rectangle has side lengths of $(x+1)$ and $(x+3)$.\\
The equation $${A(x)=x^2+4x+3}$$ gives the area of the rectangle.
%=========
To start a new paragraph we use a hard return by insert a blank line in between (like how I just do).\\
If you want to type in the same paragraph you have to use the soft return and that by ending the previous line of code with double slash (Just like how I did with the previous line).
%=========

\bigskip % to insert a blank line between (two paragraphs)
You can turn indenting off in your document by use \{\textbackslash parindent 0px\}:
%=========

\parindent 0px
This is the first sentence.\\
This is the second sentence without inserting a blank line. 
%=========

\bigskip
And you can change the distence of the indenting by increasing the px number.

\parindent 35px % The standerd is 18pt
There three ways to insert (Horizontal line)

For a simple horizontal line type \textbackslash hrule

\hrule
\bigskip
You can use \textbackslash bigskip to insert lines around the horizontal line.
Horizontal line of a specified thickness\\ % The \\ is obligatory
\rule{\linewidth}{2pt}

Controlling the horizontal line width\\
(Without \textbackslash \textbackslash)
\noindent\rule{5cm}{2pt} \\
(With \textbackslash \textbackslash)\\
\noindent\rule{5cm}{2pt} 
%=========
\end{document}
